% Part: first-order-logic
% Chapter: tableaux
% Section: proof-theoretic-notions

\documentclass[../../../include/open-logic-section]{subfiles}

\begin{document}

\iftag{FOL}
      {\olfileid[pt]{fol}{tab}{ptn}}
      {\olfileid[pt]{pl}{tab}{ptn}}
      
\olsection{Noções da teoria da prova}

\begin{editorial}
Esta seção reúne as definições da relação de demonstrabilidade e de
consistência para tableaux.
\end{editorial}

\begin{explain}
Assim como definimos várias noções semânticas importantes (validade,
consequência, satisfatibilidade), agora definimos as \emph{noções da
teoria da prova} correspondentes. Elas não são definidas recorrendo à
satisfação de !!{sentence}s em !!{structure}s, mas recorrendo à
existência de certos !!{tableau}x fechados. Foi uma descoberta
importante que essas noções coincidem. O fato de coincidirem é o
conteúdo do \emph{teorema de correção} e do \emph{teorema de completude}.
\end{explain}


\begin{defn}[Teoremas]
Uma !!{sentence}~$!A$ é um \emph{teorema} se existe um !!{tableau}
fechado para~$\sFmla{\False}{!A}$. Escrevemos $\Proves !A$ se $!A$ for
um teorema, e $\Proves/ !A$ se não for.
\end{defn}

\begin{defn}[!!^{derivability}]
!!^a{sentence} $!A$ é \emph{!!{derivable} de} um conjunto de
!!{sentence}s~$\Gamma$, $\Gamma \Proves !A$ sse existe um conjunto
finito $\{!B_1, \dots, !B_n\} \subseteq \Gamma$ e um !!{tableau}
fechado para o conjunto
\[
\{
  \sFmla{\False}{!A},
  \sFmla{\True}{!B_1}, \dots,
  \sFmla{\True}{!B_n}
  \}.
\]
Se $!A$ não for !!{derivable} de $\Gamma$, escrevemos $\Gamma \Proves/
!A$.
\end{defn}

\begin{defn}[Consistência]
Um conjunto de !!{sentence}s~$\Gamma$ é \emph{inconsistente} sse
existe um conjunto finito $\{!B_1, \dots, !B_n\} \subseteq \Gamma$ e
um !!{tableau} fechado para o conjunto
\[
\{
  \sFmla{\True}{!B_1}, \dots,
  \sFmla{\True}{!B_n}  
  \}.
\]
Se $\Gamma$ não for inconsistente, dizemos que ele é
\emph{consistente}.
\end{defn}

\begin{prop}[Reflexividade]
\ollabel{prop:reflexivity}
Se $!A \in \Gamma$, então $\Gamma \Proves !A$.
\end{prop}

\begin{proof}
Se $!A \in \Gamma$, $\{!A\}$ é um subconjunto finito de~$\Gamma$, e o
!!{tableau}
\begin{oltableau}
  [\sFmla{\False}{\formula{A}}, just = \TAss
    [\sFmla{\True}{\formula{A}}, just = \TAss,close]
  ]
\end{oltableau}
é fechado.
\end{proof}
  

\begin{prop}[Monotonicidade]
\ollabel{prop:monotonicity}
Se $\Gamma \subseteq \Delta$ e $\Gamma \Proves !A$, então $\Delta
\Proves !A$.
\end{prop}

\begin{proof}
Todo subconjunto finito de~$\Gamma$ também é um subconjunto finito
de~$\Delta$.
\end{proof}

\begin{prop}[Transitividade]
\ollabel{prop:transitivity}
Se $\Gamma \Proves !A$ e $\{!A\} \cup
\Delta \Proves !B$, então $\Gamma \cup \Delta \Proves !B$.
\end{prop}

\begin{proof}
Se $\{!A\} \cup \Delta \Proves !B$, então existe um subconjunto finito
$\Delta_0 = \{!C_1, \dots, !C_n\} \subseteq \Delta$ tal que
\begin{align*}
\{\sFmla{\False}{!B}, & \sFmla{\True}{!A}, \sFmla{\True}{!C_1},
\dots, \sFmla{\True}{!C_n}\}
\intertext{tem um !!{tableau} fechado. Se $\Gamma \Proves !A$, então
  existem $!D_1$, \dots, $!D_m \subseteq \Gamma$ tais que}
\{\sFmla{\False}{!A}, & \sFmla{\True}{!D_1},
\dots, \sFmla{\True}{!D_m}\}
\end{align*}
tem um !!{tableau} fechado.

Agora considere o !!{tableau} com as premissas
\[
\sFmla{\False}{!B},
\sFmla{\True}{!C_1}, \dots, \sFmla{\True}{!C_n},
\sFmla{\True}{!D_1}, \dots, \sFmla{\True}{!D_m}.
\]
Aplique a regra \Cut{} a~$!A$. Isso gera dois ramos: um contém
$\sFmla{\True}{!A}$ e o outro contém $\sFmla{\False}{!A}$. Assim, no
primeiro ramo, todas as premissas
\[
\{\sFmla{\False}{!B}, \sFmla{\True}{!A},
\sFmla{\True}{!C_1}, \dots, \sFmla{\True}{!C_n}\}
\]
estão disponíveis. Como existe um !!{tableau} fechado para essas
premissas, podemos anexá-lo a esse ramo; todo ramo que passa por
$\sFmla{\True}{!A}$ se fecha. No outro ramo, todas as premissas
\[
\{\sFmla{\False}{!A}, \sFmla{\True}{!D_1}, \dots,
\sFmla{\True}{!D_m}\}
\]
estão disponíveis, de modo que também podemos completar o outro lado e
obter um !!{tableau} fechado. Isso mostra que $\Gamma \cup \Delta
\Proves !B$.
\end{proof}

Observe que isso significa, em particular, que, se $\Gamma \Proves !A$
e $!A \Proves !B$, então $\Gamma \Proves !B$. Também se segue que, se
$!A_1, \dots, !A_n \Proves !B$ e $\Gamma \Proves !A_i$ para cada~$i$,
então $\Gamma \Proves !B$.

\begin{prop}
\ollabel{prop:incons}
$\Gamma$ é inconsistente sse $\Gamma \Proves !A$ para toda
!!{sentence}~$!A$.
\end{prop}

\begin{proof}
Exercício.
\end{proof}

\tagprob{FOL}
\begin{prob}
Demonstre \olref[fol][tab][ptn]{prop:incons}
\end{prob}
\tagendprob

\tagprob{notFOL}
\begin{prob}
Demonstre \olref[pl][tab][ptn]{prop:incons}
\end{prob}
\tagendprob

\begin{prop}[Compacidade]
  \ollabel{prop:proves-compact}
  \begin{enumerate}
  \item Se $\Gamma \Proves !A$, então existe um subconjunto finito
    $\Gamma_0 \subseteq \Gamma$ tal que $\Gamma_0 \Proves !A$.
  \item Se todo subconjunto finito de~$\Gamma$ for consistente, então
    $\Gamma$ será consistente.
  \end{enumerate}
\end{prop}

\begin{proof}
  \begin{enumerate}
    \item Se $\Gamma \Proves !A$, então existe um subconjunto finito
      $\Gamma_0 = \{!B_1, \dots, !B_n\}$ e um !!{tableau} fechado para
      \[
      \{\sFmla{\False}{!A}, \sFmla{\True}{!B_1}, \dots, \sFmla{\True}{!B_n}\}
      \]
      Esse !!{tableau} também mostra que $\Gamma_0 \Proves !A$.
    \item Se $\Gamma$ for inconsistente, então, para algum subconjunto
      finito $\Gamma_0 = \{!B_1, \dots, !B_n\}$, existe um !!{tableau}
      fechado para
      \[
      \{\sFmla{\True}{!B_1}, \dots, \sFmla{\True}{!B_n}\}
      \]
      Esse !!{tableau} fechado mostra que $\Gamma_0$ é inconsistente.
  \end{enumerate}
\end{proof}

\end{document}
